\documentclass[11pt,a4paper]{article}
\usepackage[T1]{fontenc}
\usepackage{lmodern}
\usepackage{amsmath,amssymb}
\usepackage{booktabs}
\usepackage[margin=2.5cm]{geometry}
\usepackage{hyperref}
\title{PTUM Library-Final Supplementary Note}
\author{Erdem Esa\\Independent Researcher, Istanbul, Turkey}
\date{}
\begin{document}
\maketitle
\section*{Purpose}
This note accompanies the library-final PTUM archive. It replaces an obsolete supplementary file that belonged to a different project. The main manuscript is now self-contained: implementation details and consistency figures are placed in the appendix.
\section*{Repository status}
The current archive contains manuscript source, compiled PDF, figures, journal class files, and library notes. A public \texttt{CLASS}/\texttt{Cobaya} repository should be attached before journal submission if the benchmark numerical claims are to be presented as fully reproducible pipeline results.
\section*{Safe claim boundary}
PTUM is an effective phenomenological cosmological framework. It is not presented as a UV completion, a solution of the cosmological-constant hierarchy, or a proof that dark matter and dark energy have a unique microscopic origin.
\section*{Future reproducibility target}
A later submission-ready archive should include: modified \texttt{CLASS} source, \texttt{Cobaya} YAML files, chains or compressed posterior samples, exact plotting scripts, random seeds, likelihood versions, and a one-command Docker reproduction path.
\end{document}
